\documentclass[11pt]{article}
\input{style-pc}
\usepackage{array}

\begin{document}

\entetesujet{Sujet bac 2017 -- Physique-Chimie -- Série D}

% ============================ CHIMIE ============================
\matiere{CHIMIE}{8 points}

\partie{A : vérification des connaissances}

\noindent\textbf{1. Questions à choix multiples.} Choisis la bonne réponse parmi les affirmations suivantes. Exemple : $1.5 = a\cdot5$.

\noindent\textbf{1.1.} La longueur d'onde d'une raie est donnée par :
\begin{enumerate}[label=$a\cdot$\arabic*.]
  \item $\lambda = \dfrac{E_n - E_p}{h\cdot c}$
  \item $\lambda = \dfrac{h(E_n - E_p)}{c}$
  \item $\lambda = \dfrac{h\cdot c}{E_n - E_p}$
\end{enumerate}

\noindent\textbf{1.2.} Le temps de demi-réaction d'une réaction d'ordre 1 est donné par :
\begin{enumerate}[label=$b\cdot$\arabic*.]
  \item $t = \dfrac{k}{\ln 2}$
  \item $t = \dfrac{1}{k\cdot C_0}$
  \item $t = \dfrac{\ln 2}{k}$
\end{enumerate}

\noindent\textbf{1.3.} Entre deux acides faibles, le plus fort est celui qui a une constante d'acidité :
\begin{enumerate}[label=$c\cdot$\arabic*.]
  \item plus faible
  \item plus grande
  \item nulle
\end{enumerate}

\noindent\textbf{1.4.} Le rendement d'estérification d'un alcool tertiaire pour un mélange équimolaire est :
\begin{enumerate}[label=$d\cdot$\arabic*.]
  \item 67 \% ;
  \item 5 \% ;
  \item 60 \%.
\end{enumerate}

\medskip
\noindent\textbf{2. Appariement.} Relie chaque élément-question de la colonne A à un élément-réponse de la colonne B. Exemple : $A_5 = B_7$.

\begin{center}\renewcommand{\arraystretch}{1.6}
\begin{tabular}{|p{0.42\linewidth}|p{0.42\linewidth}|}
\hline
\textbf{Colonne A} & \textbf{Colonne B} \\
\hline
$A_1$ : Radioactivité $\alpha$ & $B_1$ : Réaction totale \\
\hline
$A_2$ : Radioactivité $\beta^+$ & $B_2$ : Excès de nucléons \\
\hline
$A_3$ : Réaction de saponification & $B_3$ : Excès de protons \\
\hline
$A_4$ : Hydrolyse & $B_4$ : Réaction réversible \\
\hline
\end{tabular}
\end{center}

\partie{B : application des connaissances}
On prépare un ester à odeur de rhum présent dans les boissons alcoolisées en mélangeant dans un ballon 0,40 mol d'acide méthanoïque ($\ch{HCOOH}$) et 1,00 mol d'éthanol ($\ch{CH_3-CH_2-OH}$). On ajoute quelques gouttes d'acide sulfurique puis on chauffe à reflux pendant 4 heures. Après refroidissement, on dose l'acide méthanoïque présent dans le ballon par une solution d'hydroxyde de sodium $(\ch{Na^+} + \ch{OH^-})$ de concentration molaire $C_b = 1{,}6$ mol$\cdot$L$^{-1}$. Le volume de base versé pour doser tout l'acide méthanoïque restant est $V_b = 30$ mL.

\begin{enumerate}
  \item Écris l'équation-bilan de la réaction d'estérification qui a lieu puis nomme l'ester formé.
  \item Écris l'équation-bilan de la réaction de dosage de l'acide méthanoïque par la base.
  \item En te servant de la réaction de dosage, détermine (en mol) la quantité d'acide méthanoïque présent à l'équilibre.
  \item Déduis la composition (en mol) du mélange final.
  \item Calcule le rendement de la réaction.
\end{enumerate}

% ============================ PHYSIQUE ============================
\matiere{PHYSIQUE}{12 points}

\partie{A : vérification des connaissances}

\noindent\textbf{1. Questions à alternative vrai ou faux.} Réponds par vrai ou faux aux affirmations suivantes. Exemple : f = Vrai.
\begin{enumerate}[label=\alph*.]
  \item La période de rotation de la Terre est $T = 24$ h.
  \item Un ventre de vibration est un point qui vibre avec une amplitude nulle.
  \item La distance parcourue par une onde pendant une période est appelée longueur d'onde.
  \item L'allure de la trajectoire du projectile dépend de sa masse.
\end{enumerate}

\medskip
\noindent\textbf{2. Texte à trous.} Complète les mots manquants dans la phrase suivante par les mots ci-après : \textit{rectiligne, point, galiléen, mouvement, centre, isolé}.

\textit{Dans le référentiel $\cdots\cdots$, le mouvement du $\cdots\cdots$ d'inertie d'un solide $\cdots\cdots$ est un mouvement $\cdots\cdots$ uniforme.}

\medskip
\noindent\textbf{3. Question à réponse courte.} Définis l'interfrange.

\partie{B : application des connaissances}

\begin{minipage}{0.6\linewidth}
Un pendule est constitué d'une bille de masse $m$, assimilable à un solide ponctuel, fixée à une extrémité d'une tige indéformable, sans masse et de longueur $l$. L'autre extrémité de la tige peut tourner autour d'un axe horizontal $(\Delta)$ (voir figure).

\begin{enumerate}
  \item Fais le bilan des forces exercées sur la bille.
  \item Calcule le travail de ces forces lorsque le pendule, écarté d'un angle $\alpha_m = 60^\circ$ par rapport à la verticale, puis abandonné à lui-même, repasse par sa position d'équilibre.
\end{enumerate}
\end{minipage}\hfill
\begin{minipage}{0.36\linewidth}
\centering
\begin{tikzpicture}[>=Stealth,scale=1]
  \coordinate (O) at (0,0);
  \coordinate (I) at (0,-3.2);
  \coordinate (S) at ({3.2*sin(58)},{-3.2*cos(58)});
  \draw[dashed] (O)--(I);
  \draw[thick] (O)--(S);
  % angle
  \draw (0,-0.8) arc[start angle=270,end angle={270+58},radius=0.8];
  \node at (0.45,-0.75){$\alpha_m$};
  % right angle
  \draw (0.2,0)--(0.2,-0.2)--(0,-0.2);
  % arc trajectoire
  \draw[dashed] (S) arc[start angle={-90+58},end angle=-90,radius=3.2];
  \node[right] at (S){$(m)$};
  \node[right] at (1.6,-1.4){$l$};
  \node[below] at (I){$I$};
\end{tikzpicture}
\end{minipage}

\begin{enumerate}[start=3]
  \item Exprime en fonction de $m$, $l$, $g$ et $\alpha_m$, la vitesse $v$ du pendule au passage à la verticale, la vitesse initiale étant nulle.
  \item Détermine la tension de la tige au passage par la verticale.
\end{enumerate}

\partie{C : résolution d'un problème}
On réalise l'expérience des interférences lumineuses avec le dispositif des fentes de Young. La distance entre la source $S$ monochromatique et le plan des fentes $S_1$ et $S_2$ est $d = 50$ cm et la distance entre les fentes est $a = 3$ mm. L'écran d'observation $(E)$ est placé à la distance $D = 2$ m du plan des fentes (voir figure).

\begin{center}
\begin{tikzpicture}[>=Stealth,scale=1]
  \draw[dashed] (-0.3,0)--(7.6,0);
  \fill (0,0) circle (1.5pt) node[left]{$S$};
  \draw[thick] (3,-1.6)--(3,-0.5);
  \draw[thick] (3,0.5)--(3,1.6);
  \node at (2.7,0.9){$S_1$};
  \node at (2.7,-0.9){$S_2$};
  \draw[<->] (3.25,-0.45)--(3.25,0.45) node[midway,right]{$a$};
  \draw[thick] (7.4,-1.8)--(7.4,1.8) node[above right]{$(E)$};
  \draw[<->] (0,-1.4)--(3,-1.4) node[midway,below]{$d$};
  \draw[<->] (3,-1.4)--(7.4,-1.4) node[midway,below]{$D$};
\end{tikzpicture}
\end{center}

\noindent Au cours de cette expérience, on veut déterminer l'épaisseur $e$ d'une lame de verre d'indice de réfraction $n = 1{,}5$. Pour cela, on mesure sur l'écran $(E)$ la distance entre la 6\textsuperscript{e} frange brillante située d'un côté de la frange centrale et la 6\textsuperscript{e} frange brillante située de l'autre côté de la frange centrale, on trouve $L = 4{,}8$ mm.

\begin{enumerate}
  \item Détermine la longueur d'onde $\lambda_1$ de la lumière émise par la source $S$.
  \item On déplace la source $S$ parallèlement au plan des fentes, du côté de $S_1$, de $Y = 2{,}5$ cm. On constate un déplacement vertical $x$ du système de franges sur l'écran.
  \begin{enumerate}[label=\alph*.]
    \item Établis l'expression de la différence de marche $\delta$ en fonction de $Y$, $x$, $D$, $d$ et $a$.
    \item De combien et dans quel sens se déplace la frange centrale ?
  \end{enumerate}
  \item On utilise une lame de verre pour ramener la frange centrale à sa position initiale.
  \begin{enumerate}[label=\alph*.]
    \item Devant quelle fente doit-on placer la lame ?
    \item Détermine l'épaisseur $e$ de la lame.
  \end{enumerate}
\end{enumerate}

\end{document}
